% Sec. VII Replay: the recordings as a test bench (about 0.6 page).
% Paragraphs: model; beyond the link (the other recorded conditions);
% interface; positioning.
\section{Replay: The Recordings as a Test Bench}
\label{sec:replay}
This section answers whether any collaborative method can be run on
what the agents sensed and against the link they actually had.

\paragraph{Model.} A message of size $s$ sent at time $t$ over link
$(a,b)$ arrives at
\begin{equation}
  t + d_{ab}(t) + s / G_{ab}(t),
  \label{eq:replay}
\end{equation}
where $d_{ab}(t)$ is the one-way delay and $G_{ab}(t)$ the most recent
goodput measured on that link, and is dropped when the link was down
(\texttt{associated=false}, or a run of lost pings). Because the
measured clock offset is below \SI{1}{\milli\second}, $d_{ab}$ can be
measured directly from stamped probes where they exist, instead of
being assumed to be $\mathrm{RTT}/2$; only a synchronized testbed can
do this. The strict \SI{12}{\milli\second} chrony bound exceeds the
idle one-way delay, so one-way-delay claims rest on the measured
offset, and where probes are absent the module falls back to
$\mathrm{RTT}/2$ and says so. The ideal channel (instant, lossless), the synthetic channels
of prior benchmarks (a fixed bandwidth budget and a fixed or sampled
delay \cite{where2comm,v2xvit,lei2022syncnet,wei2023cobevflow}), and
the logged channel all pass through one module
(Fig.~\ref{fig:replay}), so the channel is the only variable between
runs of a method.

\paragraph{Beyond the link.} The link is one of the conditions the
recordings carry; the others are replayed the same way, each in the
same modes. \emph{Sensors}: a frame that was not recorded is missing in
replay, and the agent is told so (logged); every frame is present
(ideal, where the twin supplies it); or frames drop with each sensor's
measured rate and burst length (fitted). \emph{Pose}: each agent uses
the reference (ideal), its own recorded estimate (logged), or the
reference perturbed with the measured estimate error (fitted). A
method is thus scored with each real-world condition on or off, one at
a time, and a failure is attributed to the condition that caused it
(Sec.~\ref{sec:gym}).

\paragraph{Interface.} Any method that emits (sender, receiver, $t$,
size) plugs in and receives, for each message, an arrival time or a
drop. The implementation extends an OpenCOOD-style wrapper that already
turns latency into staleness and drops collaborators, with a schedule
driven by the trace; a converter places recorded frames in the layout
it reads. \todo{Name the released modules.}

\paragraph{Positioning.} Trace-driven network emulation replays a link
for network software, from trace modulation \cite{noble1997trace}
through netem \cite{hemminger2005netem}, Cellsim
\cite{winstein2013sprout}, Mahimahi \cite{netravali2015mahimahi}, and
Pantheon \cite{yan2018pantheon}; robotics and network co-simulators
couple a robot simulator with a simulated network
\cite{calvofullana2021rosnetsim,sommer2011veins,xu2021opencda};
collaborative-perception benchmarks simulate a link for perception,
and CooperScene \cite{cooperscene} evaluates detectors under one delay
derived from mean measured throughput. SCooP-iT replays the measured
per-message link of each node pair into perception, on the same clock
as the sensors. The fluid model of Eq.~\eqref{eq:replay} is
chosen because the iperf windows give goodput, not per-packet delivery
opportunities as in Mahimahi; E3 measures how much accuracy the simpler
model loses.

\begin{figure}[t]
\centering
\fbox{\parbox[c][4cm][c]{0.95\columnwidth}{\centering
\todo{Fig.: channel-in-the-loop replay. A method emits (sender, receiver,
$t$, size); the module returns an arrival time or a drop from the ideal,
synthetic, or logged channel. Logged: $d_{ab}(t) + s/G_{ab}(t)$, dropped
when the link was down.}}}
\caption{Channel-in-the-loop replay. One module applies the ideal,
synthetic, and logged channels, so the channel is the only variable
between runs of a method.}
\label{fig:replay}
\end{figure}

